\section{Results}
\label{sec:results}

\subsection{METIS Imbalance Factor ($\ufactor$)}

\begin{table}[h]
\centering
\caption{Effect of $\ufactor$ on partisan and compactness outcomes.
         $\Delta D$ is the change from the baseline run.
         $\Delta \overline{\PP}$ is the change in mean Polsby-Popper.}
\label{tab:ufactor-results}
\begin{tabular}{ccccc}
\toprule
$\ufactor$ & $D_{\rm nat}$ & $\Delta D$ & $\overline{\PP}$ & $\Delta\overline{\PP}$ \\
\midrule
0.1\% & 211.2 & $+0.2$ & 0.318 & $+0.009$ \\
0.5\% & 211.1 & $+0.1$ & 0.314 & $+0.005$ \\
1.0\% (default) & 211.0 & 0.0 & 0.309 & 0.0 \\
2.0\% & 210.9 & $-0.1$ & 0.305 & $-0.004$ \\
5.0\% & 210.8 & $-0.2$ & 0.299 & $-0.010$ \\
\bottomrule
\end{tabular}
\end{table}

\textbf{Partisan}: The full range of $\ufactor$ from 0.1\% to 5\% shifts
$D_{\rm nat}$ by only 0.4 seats nationally---below the level of statistical
noise ($\sigma \approx 0.8$ from B.7 seed variance).
\textbf{Compactness}: $\overline{\PP}$ varies by $\pm 0.010$ around the
baseline, a 3.2\% relative change.
Tighter imbalance (lower $\ufactor$) produces slightly more compact districts.

\subsection{County Edge Weight ($\acounty$)}

\begin{table}[h]
\centering
\caption{Effect of $\acounty$ on partisan and county-split outcomes.}
\label{tab:acounty-results}
\begin{tabular}{cccc}
\toprule
$\acounty$ & $D_{\rm nat}$ & $\Delta D$ & County splits $C$ \\
\midrule
1.0  & 211.1 & $+0.1$ & 312 \\
2.0  & 211.0 & 0.0    & 287 \\
3.0 (default) & 211.0 & 0.0 & 271 \\
5.0  & 210.9 & $-0.1$ & 248 \\
10.0 & 210.9 & $-0.1$ & 201 \\
\bottomrule
\end{tabular}
\end{table}

County edge weight has a strong effect on county integrity (splits decrease
from 312 at $\acounty = 1$ to 201 at $\acounty = 10$) with negligible
partisan effect ($\Delta D \leq 0.2$ nationally).

\subsection{ConvergenceSweep Threshold ($T$)}

\begin{table}[h]
\centering
\caption{Effect of $T$ on outcomes and runtime.}
\label{tab:T-results}
\begin{tabular}{cccc}
\toprule
$T$ & $D_{\rm nat}$ & $\Delta D$ & Wall time (50 states) \\
\midrule
100  & 211.1 & $+0.1$ & 14 min \\
200  & 211.0 & 0.0    & 22 min \\
400  & 211.0 & 0.0    & 38 min \\
600 (default) & 211.0 & 0.0 & 54 min \\
1000 & 211.0 & 0.0    & 88 min \\
\bottomrule
\end{tabular}
\end{table}

$T = 100$ produces a slightly higher $D_{\rm nat}$ (0.1 seats) because
it occasionally misses the Georgia global minimum (B.16 tail = 511).
For $T \geq 200$, outcomes are identical across all 50 states.
Runtime scales linearly in $T$.

\subsection{Area Swing ($\aswing$)}

\begin{table}[h]
\centering
\caption{Effect of $\aswing$ on partisan and compactness outcomes.}
\label{tab:aswing-results}
\begin{tabular}{cccc}
\toprule
$\aswing$ & $D_{\rm nat}$ & $\Delta D$ & $\overline{\PP}$ \\
\midrule
1.05 & 211.1 & $+0.1$ & 0.311 \\
1.07 & 211.1 & $+0.1$ & 0.310 \\
1.10 (default) & 211.0 & 0.0 & 0.309 \\
1.15 & 210.9 & $-0.1$ & 0.308 \\
1.20 & 210.9 & $-0.1$ & 0.307 \\
\bottomrule
\end{tabular}
\end{table}

\subsection{VRA Boost Weight ($\wvra$)}

\begin{table}[h]
\centering
\caption{Effect of $\wvra$ on partisan and compactness outcomes.}
\label{tab:wvra-results}
\begin{tabular}{cccc}
\toprule
$\wvra$ & $D_{\rm nat}$ & $\Delta D$ & $\overline{\PP}$ \\
\midrule
0.2 & 211.1 & $+0.1$ & 0.311 \\
0.3 & 211.0 &  0.0   & 0.310 \\
0.4 (default) & 211.0 & 0.0 & 0.309 \\
0.5 & 211.0 &  0.0   & 0.308 \\
0.6 & 210.9 & $-0.1$ & 0.306 \\
\bottomrule
\end{tabular}
\end{table}

Higher VRA boost slightly reduces compactness (as expected: minority
communities are protected at some compactness cost) with negligible
partisan effect.

\subsection{Joint Parameter Sweep ($\ufactor \times \acounty$)}
\label{sec:joint-sweep}

To test whether interactions between parameters amplify partisan effects
beyond the single-parameter bounds, we ran a full factorial sweep of
$\ufactor \times \acounty$ (3 levels $\times$ 4 levels = 12 combinations)
for Wisconsin and North Carolina, the two states with the most contested
partisan outcomes in prior papers.

\begin{table}[h]
\centering
\caption{Joint factorial sweep: $D$-seat deviation from baseline for
         Wisconsin (WI) and North Carolina (NC) across all 12 $\ufactor
         \times \acounty$ combinations.  Maximum deviation shown.}
\label{tab:joint-sweep}
\begin{tabular}{cccc}
\toprule
$\ufactor$ & $\acounty$ & $\Delta D$ (WI) & $\Delta D$ (NC) \\
\midrule
0.5\% & 1.0  & $+0.3$ & $+0.2$ \\
0.5\% & 3.0  & $+0.2$ & $+0.1$ \\
0.5\% & 5.0  & $+0.2$ & $+0.1$ \\
0.5\% & 10.0 & $+0.1$ & $+0.1$ \\
1.0\% & 1.0  & $+0.1$ & $+0.1$ \\
1.0\% & 3.0  & $0.0$  & $0.0$  \\
1.0\% & 5.0  & $-0.1$ & $0.0$  \\
1.0\% & 10.0 & $-0.1$ & $-0.1$ \\
2.0\% & 1.0  & $0.0$  & $+0.1$ \\
2.0\% & 3.0  & $-0.1$ & $-0.1$ \\
2.0\% & 5.0  & $-0.2$ & $-0.2$ \\
2.0\% & 10.0 & $-0.4$ & $-0.3$ \\
\bottomrule
\end{tabular}
\end{table}

Maximum partisan deviation in the joint sweep: $\Delta D$-seats $\leq 0.4$
(WI) and $\leq 0.3$ (NC).
The interaction effect does not amplify beyond the single-parameter
bounds; the worst-case cell ($\ufactor = 2.0\%$, $\acounty = 10$) is
consistent with the additive upper bound from Section~\ref{sec:key-finding}.
